\section{Future Directions}
\begin{frame}{Future Directions in Attention Mechanisms}
  \small
  \begin{itemize}
    \item[$\blacktriangleright$] \textbf{Approximate and Sparse Attention:}
      \begin{itemize}
        \item \textit{Linformer, Performer, Longformer, Routing Transformer} -- Replace the dense score matrix with a low-rank, kernelized or sparse stand-in, trading exactness for sub-quadratic cost.
      \end{itemize}
    \item[$\blacktriangleright$] \textbf{Exact but IO-Aware Attention:}
      \begin{itemize}
        \item \textit{FlashAttention} -- Computes the same attention, tiled to cut traffic to GPU memory. It saves wall-clock time rather than arithmetic.
      \end{itemize}
    \item[$\blacktriangleright$] \textbf{Sparse Normalization:}
      \begin{itemize}
        \item \textit{Sparsemax} -- A softmax replacement that can assign exactly zero weight, giving sharper and more interpretable attention maps.
      \end{itemize}
    \item[$\blacktriangleright$] \textbf{Cross-modal Attention:}
      \begin{itemize}
        \item \textit{Flamingo, BLIP-2} -- Cross-attention layers let a language model read the output of a vision encoder.
      \end{itemize}
    \item[$\blacktriangleright$] \textbf{Learnable Routing:}
      \begin{itemize}
        \item \textit{Mixture-of-Experts, Mixture-of-Depths} -- A learned gate sends each token to a subset of the parameters or layers.
      \end{itemize}
  \end{itemize}
\end{frame}
